\subsubsection{Benchmarking Attention-Based Interpretability of Deep Learning in Multivariate Time Series Predictions}

\citet{baric2021benchmarking} introduce ten synthetic multivariate time-series
systems with progressively more complex dynamics. The collection includes
independent autoregressive series, nonlinear autoregressive equations,
cross-series interactions, switching rules, a two-dimensional Ising system, and
a logistic-map-inspired system. For a target series $x_{k,t}$, the known
generating equation defines a \textbf{binary feature--lag mask}: each entry is
1 when the corresponding lagged source occurs in the equation and 0 when it
does not. This structural ground-truth explanation can be represented as

\begin{equation}
    M^{\mathrm{GT}}_{k,j,\ell}
    =
    \mathbb{I}\!\left[
        x_{j,t-\ell}\text{ occurs in the generating equation for }x_{k,t}
    \right],
\end{equation}

where $j$ identifies the source series and $\ell$ identifies its lag. For
example, an autoregressive equation containing lags 3 and 7 defines those two
lags as the expected temporal explanation. Cross-series equations similarly
define the expected source-variable explanation. This equation is anchored at
the target time $t$: lag 1 refers to $t-1$, the latest input available for that
target. Equivalently, setting the forecast origin to $\tau=t-1$ places that
sample at relative time 0 and the target $x_{k,t}$ at relative time $+1$.

\begin{landscape}
\begingroup
\scriptsize
\setlength{\LTpre}{0pt}
\setlength{\LTpost}{0pt}
\setlength{\LTcapwidth}{0.92\linewidth}
\setlength{\tabcolsep}{4pt}
\renewcommand{\arraystretch}{0.94}
\begin{longtable}{@{}
>{\raggedright\arraybackslash}p{0.10\linewidth}
>{\raggedright\arraybackslash}p{0.46\linewidth}
>{\raggedright\arraybackslash}p{0.38\linewidth}@{}}
\caption{Exact dataset definitions from Table~1 of
\citet{baric2021benchmarking}: models used for benchmarking.}
\label{tab:baric-datasets-exact}\\
\toprule
Name & Formula & Parameters \\
\midrule
\endfirsthead
\multicolumn{3}{@{}l}{\small\emph{Table~\thetable\ continued}}\\
\toprule
Name & Formula & Parameters \\
\midrule
\endhead
\midrule
\multicolumn{3}{r@{}}{\small\emph{Continued on the next page}}\\
\endfoot
\bottomrule
\endlastfoot

Dataset 1 &
$X_{n,t}=C_n+\varepsilon_t$ &
$C_n=\mathcal{N}(0,1)$ \\

Dataset 2 &
$X_{n,t}=c_{t_{\mathrm{lag}}}X_{n,t-t_{\mathrm{lag}}}+\varepsilon_t$ &
$c_3=1/2,\ c_7=1/2$ \\

Dataset 3 &
$X_{n,t}=\tanh\!\left(c_{t_{\mathrm{lag}}}X_{n,t-t_{\mathrm{lag}}}+\varepsilon_t\right)$ &
$c_3=5/7,\ c_7=1/7,\ c_9=1/7$ \\

Dataset 4 &
$X_{1-n,t}=c_{t_{\mathrm{lag}}}X_{n,t-t_{\mathrm{lag}}}+\varepsilon_t$ &
$c_2=2/5,\ c_5=1/5,\ c_9=2/5$ \\

Dataset 5 &
$X_{n,t}=c_{n,t_{\mathrm{lag}}}X_{0,t-t_{\mathrm{lag}}}+\varepsilon_t$ &
$\begin{aligned}
c_{0,3}&=1/2,&c_{0,4}&=1/2,\\
c_{1,9}&=1,\\
c_{2,2}&=1/2,&c_{2,7}&=1/2,\\
c_{3,3}&=1/10,&c_{3,4}&=1/10,&c_{3,8}&=4/5,\\
c_{4,2}&=1/3,&c_{4,5}&=2/9,&c_{3,8}&=4/9
\end{aligned}$ \\

Dataset 6 &
$X_{n,t}=\tanh\!\left(C_{n,t_{\mathrm{lag}}}X_{0,t-t_{\mathrm{lag}}}+\varepsilon_t\right)$ &
$\begin{aligned}
c_{0,3}&=1/2,&c_{0,4}&=1/2,\\
c_{1,9}&=1,\\
c_{2,2}&=1/2,&c_{2,7}&=1/2,\\
c_{3,3}&=1/10,&c_{3,4}&=1/10,&c_{3,8}&=4/5,\\
c_{4,2}&=1/3,&c_{4,5}&=2/9,&c_{3,8}&=4/9
\end{aligned}$ \\

Dataset 7 &
$\begin{aligned}
X_{0,t}&=c_{0,1}X_{0,t-1}+c_{0,5}X_{0,t-5}+\varepsilon_t,\\
X_{1,t}&=1+c_{1,2}X_{0,t-2}+\varepsilon_t,\\
X_{2,t}&=c_{2,1}X_{1,t-1}+c_{2,4}X_{3,t-4}+\varepsilon_t,\\
X_{3,t}&=1+c_{3,4}X_{2,t-4}+c_{3,1}X_{0,t-1}+\varepsilon_t,\\
X_{4,t}&=c_{4,4}X_{4,t-4}+c_{4,1}X_{1,t-1}+\varepsilon_t
\end{aligned}$ &
$\begin{aligned}
c_{0,1}&=1/4,&c_{0,5}&=3/4,\\
c_{1,2}&=-1,\\
c_{2,1}&=1,&c_{2,4}&=1,\\
c_{3,4}&=-2/7,&c_{3,1}&=5/7,\\
c_{4,4}&=12/22,&c_{4,1}&=10/22
\end{aligned}$ \\

Dataset 8 &
$\begin{aligned}
&\text{if }X_{0,t-5}>1/2:\\
&X_{0,t}=c_{0,1}X_{0,t-1}+c_{0,3}X_{0,t-3}+\varepsilon_t,\\
&X_{1,t}=X_{0,t-5}+\varepsilon_t,\\
&X_{2,t}=X_{0,t-4}+\varepsilon_t,\\
&X_{3,t}=c_{3,1}X_{3,t-1}+c_{3,4}X_{3,t-4}+\varepsilon_t,\\
&\text{else:}\\
&X_{0,t}=c_{0,1}X_{0,t-1}+c_{0,3}X_{0,t-3}+\varepsilon_t,\\
&X_{1,t}=X_{3,t-2}+\varepsilon_t,\\
&X_{2,t}=X_{3,t-4}+\varepsilon_t,\\
&X_{3,t}=c_{3,1}X_{3,t-1}+c_{3,4}X_{3,t-4}+\varepsilon_t
\end{aligned}$ &
$\begin{aligned}
c_{0,1}&=1/2,&c_{0,3}&=1/2,\\
c_{3,1}&=1/2,&c_{3,4}&=1/2
\end{aligned}$ \\

Dataset 9 &
$H(\sigma)=-\sum_{\langle i,j\rangle}J_{i,j}\sigma_i\sigma_j
-\mu\sum_j h_j\sigma_j$ &
$T=2,\ T_c,\ 2.75$ \\

Dataset 10 &
$\begin{aligned}
X_{0,t}&=rX_{0,t-3}(1-X_{0,t-3}),\\
X_{1,t}&=rX_{1,t-5}(1-X_{1,t-5}),\\
X_{2,t}&=\tfrac12X_{0,t-3}+\tfrac12X_{1,t-5}
\end{aligned}$ &
$r=1.5,\ 2.5,\ 3.2,\ 3.55,\ 3.56996$ \\

\end{longtable}
\endgroup
\end{landscape}

\paragraph{Reading the original synthetic-dataset table.}
Table~\ref{tab:baric-datasets-exact} above transcribes the formulas and
parameters in Table~1 of \citet{baric2021benchmarking}. Mathematical typography
and line breaking have been adapted to this document; dataset names, equations,
and parameter values are retained from the source table. The source caption
names the logistic-map system as Dataset~9 and the Ising system as Dataset~10,
while the table rows and the paper's later experimental discussion associate
the Ising Hamiltonian with Dataset~9 and the logistic-map equations with
Dataset~10. The row definitions preserve the operational specification shown
by the authors.

\paragraph{How each model produces an explanation.}
The benchmark reads the explanations from internal attention or
dependency-discovery quantities learned during forecasting. The four outputs
have different meanings and therefore require model-specific interpretation.

\textbf{seq2graph.}
In seq2graph \citep{dang2018seq2graph}, a dual-purpose recurrent unit first
processes every component series separately. Its temporal attention produces
$\alpha^{\mathrm{att}}$ coefficients over the $L$ positions of the input
window. For source series $j$, $\alpha^{\mathrm{att}}_{j,\ell}$ indicates how much the representation attends to
lag $\ell$, with the temporal weights summing to one. The lag-weighted
representations of all series are then concatenated and passed to a decoder.
The decoder produces $\beta^{\mathrm{att}}_{k,j}$ coefficients over source
series for target $k$, again summing to one. Thus, the benchmark interprets
$\alpha^{\mathrm{att}}$ as a soft lag or autocorrelation explanation and
$\beta^{\mathrm{att}}$ as a directed source-variable or cross-series
explanation.

\textbf{IMV-LSTM.}
IMV-LSTM \citep{guo2019imvlstm} preserves a hidden-state matrix whose row $j$
contains the representation of input variable $j$. Temporal attention over
this matrix yields $\alpha^{\mathrm{att}}_{j,\ell}$ and therefore identifies the historical
positions emphasized inside each variable representation. The model combines
the temporally weighted hidden state with the original hidden state, and a
mixture-attention layer then produces $\beta^{\mathrm{att}}_{k,j}$, the
importance assigned to source variable $j$
when forecasting target $k$. A separate target-specific model is trained in
the benchmark, and stacking its $\beta^{\mathrm{att}}$ vectors produces a
target-by-source importance matrix. Together, $\alpha^{\mathrm{att}}$
describes when the model attends and $\beta^{\mathrm{att}}$ describes which
variable it uses.

\textbf{TCDF.}
TCDF \citep{nauta2019tcdf} trains one attention-based temporal convolutional
network for every target series. A learned attention vector gates the $N$ input
channels and supplies candidate source variables. Semi-binarization converts
these soft channel weights into a selected set of potential causes. In the full
TCDF procedure, permutation-importance validation tests whether perturbing a
candidate source reduces prediction quality, and the largest relevant
depthwise-convolution kernel weight supplies its delay. The final explanation
is therefore a set of directed pairs with delays,

\begin{equation}
    \widehat{\mathcal{E}}_k
    = \{(j,\widehat{\ell}_{k,j}) :
    x_j\text{ is retrieved as a cause of }x_k\},
\end{equation}

rather than a normalized heatmap of soft attention. Bari\'c et al. summarize
this discrete output by the percentage of repeated experiments in which each
target--source relation is retrieved.

\textbf{DA-RNN.}
DA-RNN \citep{qin2017darnn} uses two attention stages. At encoder step $t$, the
input-attention layer assigns weights to the driving series and multiplies each
current input by its weight before the LSTM update. Aggregating these weights
over time and examples produces a source-variable score comparable to a
$\beta^{\mathrm{att}}$ matrix. In the decoder, temporal attention weights the
sequence of encoder hidden states. Aggregating these weights produces an
$\alpha^{\mathrm{att}}$-like
temporal score. These temporal weights index encoder representations, each of
which already combines the selected driving variables; consequently, they
provide a time-only explanation rather than a source-specific feature--lag
map. The target history also enters outside the driving-series attention, so
the variable-attention matrix lacks a directly comparable self-dependency entry
on its diagonal.

\paragraph{How the outputs are compared with the generating formula.}
For the linear generators, write the known mechanism as

\begin{equation}
    x_{k,t}
    = b_k + \sum_{j=1}^{N}\sum_{\ell=1}^{L}
      c_{k,j,\ell}x_{j,t-\ell} + \varepsilon_{k,t}.
\end{equation}

The nonzero coefficients define two reference objects. The feature--lag mask
$M^{\mathrm{GT}}_{k,j,\ell}=\mathbb{I}[c_{k,j,\ell}\neq 0]$ specifies exactly
which source and lag occur in the equation, and the source-only adjacency map
is obtained by collapsing the lag axis, meaning that all lag entries for a
fixed target--source pair $(k,j)$ are reduced to one binary value that equals 1
when at least one of those lags is active,

\begin{equation}
    B^{\mathrm{GT}}_{k,j}
    = \mathbb{I}\!\left[\sum_{\ell=1}^{L}
      M^{\mathrm{GT}}_{k,j,\ell}>0\right].
\end{equation}

The nonlinear and switching datasets use the same occurrence rule: a source--lag
pair belongs to the reference when it participates in the active generating
equation. The paper then aligns each model output with the compatible projection
of this reference:

\begin{itemize}
    \item seq2graph and IMV-LSTM $\beta^{\mathrm{att}}$ heatmaps are compared
    with $B^{\mathrm{GT}}$. Their $\alpha^{\mathrm{att}}$ heatmaps are compared
    with the nonzero lag locations in $M^{\mathrm{GT}}$.
    \item TCDF target--source retrieval frequencies are compared with
    $B^{\mathrm{GT}}$, while its retrieved delays are checked against the
    nonzero lag indices of $M^{\mathrm{GT}}$.
    \item aggregated DA-RNN input attention is compared with
    $B^{\mathrm{GT}}$. Its temporal attention is compared only at the time axis
    because the encoder-state index cannot be resolved back to an individual
    input variable.
\end{itemize}

The expected patterns are obtained directly from the equations. In Dataset 2,
every series depends only on its own lags 3 and 7, so the expected source matrix
is the identity and the expected temporal peaks occur at lags 3 and 7. In
Dataset 4, each series is generated from the other series, so the expected
source matrix has high values on the anti-diagonal. In Dataset 5, the first
series drives the remaining series, so the expected source matrix concentrates
on its first column.

For the soft-attention models, the authors average the learned coefficients and
display target--source or target--lag heatmaps beside the equation-derived
pattern. Standard deviations across repeated training experiments indicate
whether the displayed explanation is stable. The benchmark uses five runs for
seq2graph, three for IMV-LSTM, ten for TCDF, and three for DA-RNN. For TCDF,
each heatmap cell is instead a retrieval frequency, so a value of $0.7$ means
that the corresponding directed relation was found in 70\% of runs; a row may
sum to less than one when some runs retrieve no cause.

This evaluation tests whether high importance or retrieved edges occur at the
correct support locations. Since the attention coefficients are nonnegative
and normalized, their numerical values represent relative importance rather
than the signed coefficients $c_{k,j,\ell}$ of the generator. Consequently,
positive and negative generating effects share the same relevance target. The
published comparison is primarily visual and qualitative: correctness is judged
from agreement between the learned heatmap and the expected support pattern,
while mean, standard deviation, and TCDF retrieval percentages measure
confidence and reproducibility. A single tensor-level distance, F1 score, or
rank-correlation score is outside the reported benchmark evaluation.

